\cotefile{29}
\scanpage{17}
dernière remarque aurait dû figurer avant l'exemple, elle s'applique chaque fois que le foncteur $\underline{r}$ est fidèle). \deleted{Nous pouvons également} Supposons maintenant $V$ dense dans $Z$ et de plus $Z$ \uline{normal}, alors on voit aussitôt que le foncteur $Z'\qquad Z'\;V$ est pleinement fidèle, d'où résult facilement qu'il en est de même du foncteur réalisation géométrique $Z'\qquad Z'/G$; donc dans ce cas, la catégorie $\mathrm{Rev}_{\underline{R}}(S)$ \deleted{\unsure{etd|età|est}} estéquivalente à une sous-catégorie pleine de la catégorie des revêtements \deleted{\unsure{de|du}} normaux de $S$,\inserted{étales au dessus de $U$, et} dont toute composante \deleted{domine} irréductible domine une composante irréductible de $S$ rencontrant $U$.

\uline{Exemple} 2 C'est une simple variante de l'exemple précédent: au lieu de se donner un $S$-schéma $Z$ à groupe fini d'opérateurs $G$, on se donne un systême projectif $(Z_i,G_i)_{i\;I}$ de tels objets ($I$ ensemble ordonné filtrant). \deleted{La catégorie Rev} Le systême projectif de\unsure{\deleted{s}|x|s|\deleted{s}x|\deleted{x}s} groupes $(G_i)$ définit une catégorie fibrée $C_o$ sur la catégorie associée à $I$ (les fibres de $C_o$ étant les catégories à un seul objet définies par les $G_i^o$), et on prendra $\underline{C}=\underline{\mathrm{Et}}_S\mathrm{x}C_o$, et on définit $\underline{s}$ de façon évidente. La catégorie $\mathrm{Rev}_{\underline{R}}(S)$ est maintenant équivalente à la catégorie $\mathrm{Lim}$ des catégories $\mathrm{Rev}_{\underline{R}_i}(S)$ définies séparément par les $(Z_i,G_i)$. Un cas intéressant est celui où les foncteurs transition (qui s'interprètent de façon évidente comme des foncteurs imagesinverses de revêtements étales à opérateurs) sont pleinement fidèles, de sorte que les $\mathrm{Rev}_{\underline{R}_i}(S)$ s'identifient (à équivalence près) à des sous-catégories pleines de $\mathrm{Rev}_{\underline{R}}(S)$ qui epuisent évidemment cette dernière. Il en sera ainsi si les morphismes de transition $(Z_i,G_i)\qquad(Z_j,G_j)$ sont tels que $G_i\quad G_j$ soit surjecti et $Z_i\quad Z_j$ induit un isomorphisme $Z_j\quad Z_i/N_{ij}\mathord{,}_{\unsure{R|\deleted{R}}}$ où $N_{ij}$ est le noyau de $G_i\quad G_j$ (en effet, il en résulte que pour tout revêtement \inserted{étale $Z'_j$}
% archive sheet number: 16

